% !TeX root = ../../Book1.tex
\section{Banach–Alaoglu 定理}
\label{b1:sec:0903}

在范数拓扑中，无穷维空间的闭单位球不紧；换用弱-*拓扑以后，对偶单位球却总是紧的。
这个结论不要求原空间完备，也不要求可分。证明的起点很具体：
一个范数不超过一的泛函，在每个固定向量上的取值都落在一个紧圆盘内。
困难在于这些向量通常不可数，不能直接用可数次抽取子列处理。

\subsection{对偶单位球}
\label{b1:subsec:090301}

记 \(B_{X^*}=\{f\in X^*:\|f\|\leq1\}\)。对每个 \(x\in X\)，令
\[
 D_x=\{z\in\mathbb K:|z|\leq\|x\|\},
 \quad P=\prod_{x\in X}D_x.
\]
这里 \(\mathbb K=\mathbb R\) 或 \(\mathbb C\)；实情形的 \(D_x\) 是闭区间，复情形是闭圆盘。
把泛函送到它的全部取值，就得到单射
\[
 \Phi:B_{X^*}\longrightarrow P,
 \quad \Phi(f)=(f(x))_{x\in X}.
\]
乘积中的一般元素只是一份取值表，不一定满足线性关系。证明将先得到整份取值表空间的紧性，
再把线性关系作为闭条件加进去。

在乘积空间 \(\prod_{i\in I}K_i\) 上，\term[chengji-tuopu]{乘积拓扑}[product topology]的基本开集只限制有限个坐标：
\[
 \bigcap_{i\in F}\pi_i^{-1}(U_i),
 \quad F\subset I\text{ 有限},\quad U_i\subset K_i\text{ 开},
\]
其中 \(\pi_i\) 是第 \(i\) 个坐标投影。因此 \(\Phi\) 把弱-*基本邻域变成像空间中的乘积基本邻域。
这正是弱-*拓扑与逐点测试之间的关系，而不是一种另加的紧性假设。

本节暂时需要一般拓扑中的紧性：每个开覆盖有有限子覆盖。
取补集可知，这等价于每个具有\term[youxian-jiao-xingzhi]{有限交性质}[finite intersection property]的闭集族都有公共点；
有限交性质是指任意有限个成员的交非空。
下面给出乘积紧性所需的证明，以免把本章关键一步推给尚未展开的拓扑附录。
乘积拓扑和定理的一般版本可参阅 Fremlin\cite[卷三，3A3I--J]{Fremlin2016Measure}。

\begin{theorem}[Tychonoff]\label{b1:thm:tychonoff}
任意一族紧拓扑空间的乘积，在乘积拓扑下仍然紧。
\end{theorem}
\begin{proof}
若有因子为空，乘积为空，结论成立。以下设各因子 \(K_i\) 非空，记 \(P=\prod_iK_i\)。
本证明使用选择公理及其等价形式 Zorn 引理。
给定 \(P\) 中具有有限交性质的闭集族 \(\mathcal F\)，
把它扩充为具有有限交性质的极大子集族 \(\mathcal U\)。
这样的扩充存在，因为按包含关系排序的一条链的并仍具有有限交性质：
有限多个成员已同时属于链中的某一个族。

极大性给出三个性质。首先，\(\mathcal U\) 对有限交封闭，
因为加入已有成员的有限交不会破坏有限交性质。其次，若 \(A\in\mathcal U\) 且 \(A\subset B\subset P\)，
则 \(B\in\mathcal U\)，理由相同。最后，对任意 \(B\subset P\)，
\(B\) 与 \(P\setminus B\) 中恰有一个属于 \(\mathcal U\)。
若 \(B\notin\mathcal U\)，加入 \(B\) 会破坏有限交性质，因而存在有限多个已有成员，
其交 \(A\) 与 \(B\) 不交；前两个性质便推出 \(P\setminus B\in\mathcal U\)。
二者不能同时属于 \(\mathcal U\)，因为空集不属于这个族。

对每个 \(i\)，闭集族
\[
 \{\overline{\pi_i(A)}:A\in\mathcal U\}
\]
具有有限交性质；任取有限多个 \(A\) 的一个公共点并投影，就能验证这一点。
由 \(K_i\) 紧，可选
\(x_i\in\bigcap_{A\in\mathcal U}\overline{\pi_i(A)}\)，并令 \(x=(x_i)_{i\in I}\)。
若 \(U_i\) 是 \(x_i\) 的开邻域，则 \(\pi_i^{-1}(U_i)\in\mathcal U\)。
否则其补集属于 \(\mathcal U\)，而该补集的投影包含在闭集 \(K_i\setminus U_i\) 中，
与 \(x_i\) 的选取矛盾。

故 \(x\) 的每个乘积基本邻域都属于 \(\mathcal U\)，与每个 \(F\in\mathcal F\) 相交。
这说明 \(x\in\overline F=F\)。于是原闭集族有公共点，\(P\) 紧。
\end{proof}

证明中有限个坐标可以同时控制，是乘积拓扑的决定性特征。
若把“有限”改成“任意多个”，得到的拓扑一般更强，上述邻域论证便不再成立。

\subsection{Banach–Alaoglu 定理}
\label{b1:subsec:090302}

\begin{theorem}[Banach--Alaoglu]\label{b1:thm:banach-alaoglu}
设 \(X\) 是赋范空间。对偶闭单位球 \(B_{X^*}\) 在弱-*拓扑 \(\sigma(X^*,X)\) 下是紧的。
\end{theorem}
\begin{proof}
每个 \(D_x\) 紧，由\cref{b1:thm:tychonoff}，\(P\) 紧。
在 \(P\) 中考虑同时满足下列关系的取值表 \(u=(u_x)_{x\in X}\)：
\[
 u_{ax+by}=au_x+bu_y
 \quad(x,y\in X,\ a,b\in\mathbb K).
\]
对固定的 \(x,y,a,b\)，等式两边的差是连续的有限坐标函数，所以该等式定义闭集。
所有等式的公共解集因而是 \(P\) 的闭子集。

每张这样的表定义一个线性泛函 \(f(x)=u_x\)，且 \(|f(x)|\leq\|x\|\)，
故 \(f\in B_{X^*}\)。反过来，\(B_{X^*}\) 中每个泛函的取值表满足全部等式。
公共解集因此恰为 \(\Phi(B_{X^*})\)，它是紧集。
由于 \(\Phi\) 及其到像上的逆都把有限测试邻域对应起来，
\(\Phi\) 是弱-*拓扑到像的同胚，所求紧性成立。
\end{proof}

\term*[banach-alaoglu-dingli]{Banach--Alaoglu 定理}[Banach--Alaoglu theorem]中的星号不可省去。
它说的是 \(X^*\) 上由 \(X\) 测试的紧性，而不是由全部 \(X^{**}\) 测试的弱紧性。
闭球与有界性的条件也各有作用：对任意 \(R\geq0\)，半径为 \(R\) 的对偶闭球弱-*紧，
从而每个范数有界、弱-*闭的子集弱-*紧；一般范数闭集却不必弱-*闭。
Fremlin\cite[卷三，3A5F]{Fremlin2016Measure}采用同一赋范空间版本，亦不要求 \(X\) 完备。

下一节要把对偶球的紧性传回 \(X\)。为此，还需知道自然嵌入的像在二次对偶球中有多大。
它未必等于整个球，但有限多个泛函不能把球中的点与这个像分开。

\begin{theorem}[Goldstine]\label{b1:thm:goldstine}
设 \(J_X:X\to X^{**}\) 是自然等距嵌入。则
\[
 \overline{J_X(B_X)}^{\,\sigma(X^{**},X^*)}=B_{X^{**}}.
\]
\end{theorem}
\begin{proof}
右边的球由闭条件 \(|x^{**}(f)|\leq\|f\|\) 描述，故弱-*闭，并包含左边的闭包。
反过来，固定 \(x^{**}\in B_{X^{**}}\)、\(f_1,\ldots,f_m\in X^*\) 和 \(\varepsilon>0\)。
令
\[
 T:X\to\mathbb K^m,\quad Tx=(f_1(x),\ldots,f_m(x)),
 \quad z=(x^{**}(f_1),\ldots,x^{**}(f_m)).
\]
我们断言 \(z\in\overline{T(B_X)}\)。若不然，有限维空间中的凸分离定理给出一个实线性泛函，
把 \(z\) 与这个闭凸集严格分离。实情形它写成 \(w\mapsto\sum_i a_iw_i\)；
复情形则写成 \(w\mapsto\operatorname{Re}\sum_i a_iw_i\)。
记 \(g=\sum_i a_if_i\)，两种情形统一给出
\[
 \operatorname{Re}x^{**}(g)>
 \sup_{x\in B_X}\operatorname{Re}g(x)=\|g\|.
\]
最后的等号来自单位球的对称性；复情形还可乘以模为一的标量。
但 \(|x^{**}(g)|\leq\|g\|\)，矛盾。
于是能选 \(x\in B_X\)，使所有 \(|f_i(x)-x^{**}(f_i)|<\varepsilon\)。
这正是 \(x^{**}\) 属于所求弱-*闭包的邻域判据。
\end{proof}

\term*[goldstine-dingli]{Goldstine 定理}[Goldstine theorem]给出的是有限测试意义的逼近，
并不声称 \(J_X(B_X)\) 在二次对偶范数下稠密。
事实上，若 \(X\) 完备，等距像 \(J_X(X)\) 已是范数闭子空间；
除非自然嵌入满射，否则它不可能在 \(X^{**}\) 中范数稠密。

\subsection{可分情形的序列版本}
\label{b1:subsec:090303}

紧性在一般拓扑下不自动提供收敛子列。下面的可分性假设使有界对偶球只需可数个测试，
由此恢复度量空间中熟悉的子列结论。

\begin{theorem}\label{b1:thm:weak-star-metrizable}
设 \(X\) 是可分赋范空间，\(R>0\)。在 \(RB_{X^*}\) 上，弱-*拓扑可由一个度量给出。
特别地，每个范数有界的 \(X^*\) 中的序列都有弱-*收敛子列，其极限仍在 \(X^*\) 中。
\end{theorem}
\begin{proof}
取 \(X\) 的可数稠密集 \(\{x_j:j\geq1\}\)，在 \(RB_{X^*}\) 上定义
\[
 d(f,g)=\sum_{j=1}^{\infty}2^{-j}\min\{1,|(f-g)(x_j)|\}.
\]
对称性及三角不等式直接来自绝对值和
\(\min\{1,a+b\}\leq\min\{1,a\}+\min\{1,b\}\)。
若 \(d(f,g)=0\)，则 \(f,g\) 在稠密集上一致，由连续性知 \(f=g\)，所以 \(d\) 是度量。

先说明它确实给出弱-*拓扑。给定一个 \(d\)-球，先使级数尾和小于半个半径，
再限制前面有限多个测试值，就得到包含在球内的弱-*邻域。
反过来，考虑有限个条件 \(|(g-f)(y_i)|<\varepsilon\)。
分别选 \(x_{j_i}\) 逼近 \(y_i\)，使 \(2R\|y_i-x_{j_i}\|<\varepsilon/2\)。
因为
\[
 |(g-f)(y_i)|\leq |(g-f)(x_{j_i})|+2R\|y_i-x_{j_i}\|,
\]
只需把有限个 \(|(g-f)(x_{j_i})|\) 控制在 \(\varepsilon/2\) 内。
这些条件又包含 \(f\) 的某个 \(d\)-球：在距离和中分别保留第 \(j_i\) 项即可。
两种邻域系统因此等价。

由 Banach--Alaoglu 定理，这个度量空间紧，故序列紧。
也可以直接看到抽取过程：先逐个在有界标量列 \((f_n(x_j))_n\) 中抽子列，
再取对角子列。上述逼近估计保证这条子列在每个 \(x\in X\) 上取值都为 Cauchy 列。
令 \(f(x)\) 为逐点极限，线性关系可以逐点取极限，且 \(|f(x)|\leq R\|x\|\)。
于是 \(f\in RB_{X^*}\)，抽出的子列弱-*收敛到 \(f\)。
\end{proof}

范数有界在这里不可去掉：稠密集上的信息要推广到整个空间，靠的是对所有泛函共同有效的 \(2R\) 估计。
这个证明也说明，抽子列时测试函数不必事先包含空间中的每一个向量。
先在一个可数稠密族上安排收敛，再用一致估计补足其余测试，是后面弱解紧性论证的常用步骤。

不可分情形确实可能没有这样的子列。令 \(I=\mathcal P(\mathbb N)\)，
取 \(X=\ell^1(I)\)，即所有至多可数个非零坐标、绝对可和的族 \(x=(x_A)_{A\in I}\)，
范数为 \(\sum_A|x_A|\)。定义
\[
 f_n(x)=\sum_{A\in I}\mathbf1_A(n)x_A.
\]
则 \(\|f_n\|\leq1\)。任取子列 \((f_{n_k})\)，令 \(A=\{n_{2j}:j\geq1\}\)，
并以第 \(A\) 个坐标的单位向量 \(e_A\) 测试。
序列 \(f_{n_k}(e_A)\) 在零与一之间交替，不收敛。
所以这一对偶单位球虽然弱-*紧，却含有无弱-*收敛子列的序列。
一般紧性与序列紧性之间的区别不能仅凭记号中的收敛箭头消除。
